\documentclass[11pt]{article}
\usepackage[margin=1in]{geometry}
\usepackage{amsmath,hyperref}
\setlength{\emergencystretch}{1em}
\title{Corrections to the Super-Resolution Project Report}
\author{Liam Cawley}
\date{October 2026}
\begin{document}
\maketitle

The original report by Sophia Tesic, Liam Cawley, Alexandra Lavacek,
Namitha John, and Dennis Lin follows this correction. Its appended notebook
contains training output, but the published evaluation has errors that prevent
the tables from establishing held-out performance.

\section*{Training and evaluation used the same split}
In the appended \texttt{train\_and\_evaluate\_model} function, the argument named
\texttt{val\_dataset} is passed to \texttt{model.fit} and then reused for scoring.
The calling cell builds it with \texttt{subset='valid'}. These scores therefore
cannot be presented as an independent evaluation on held-out DIV2K images.
The corrected wrapper requires separate training and evaluation manifests and
rejects overlapping image IDs before fitting.

\section*{The simple model's pixel range was mishandled}
The simple model uses a tanh output and is trained on images transformed by
$x/127.5-1$. The later comparison sends it unnormalized inputs, clips its output
to $[0,255]$, and casts directly to \texttt{uint8}. An output in $[-1,1]$ first
needs the inverse transform $127.5(x+1)$.

This failure is reproducible without training a network. A perfect signed-range
prediction of a synthetic RGB gradient scores $4.77054$ dB under the old cast;
inverse normalization reconstructs the reference to numerical precision.
This demonstrates that the reported $4.80$ dB is not evidence that direct image
prediction inherently fails. It does not recover the original model's true
performance, since the input preprocessing was also inconsistent.

\section*{Metrics and comparisons}
The replacement evaluator uses RGB float images in $[0,1]$, a fixed
\texttt{data\_range=1}, and SSIM with \texttt{channel\_axis=-1} and an
11-pixel window. It scores each supplied evaluation image for both the model
and the bicubic baseline and retains per-image results.
It avoids inferring the metric range from each image's minimum and maximum.

The report contains several comparisons on different subsets and configurations.
For example, its summary table and its later side-by-side comparisons use
different bicubic averages. They cannot be combined into one estimate of gain.
Architecture, optimization, and preprocessing also changed together, so the
tables do not isolate the effect of residual learning or perceptual loss.

\section*{Available correction and remaining work}
\texttt{evaluation.py} repairs the split and preprocessing contracts and includes
the reproducible range-error check. The archived checkpoints are unavailable;
there is no new held-out model-performance claim in this correction. A valid
comparison requires training on the training split and scoring all methods on
the same untouched evaluation images under the declared pixel convention.

The implemented networks are direct image-to-image models. The diffusion
terminology describes their architectural inspiration; they do not implement a
noise-conditioned reverse-diffusion sampling process.

\noindent Code and check results:
\url{https://cawley.dev/research/super-resolution-ddpm/evaluation.py}\\
\url{https://cawley.dev/research/super-resolution-ddpm/evaluation_checks.json}
\end{document}
